\documentclass[conference]{IEEEtran}

\usepackage[T5]{fontenc}
\usepackage[utf8]{inputenc}
\usepackage{amsmath,amssymb}
\usepackage{booktabs}
\usepackage{cite}
\usepackage{graphicx}
\usepackage{svg}
\usepackage{url}

\title{TAD-YOLO: A Tiny-Aware Dynamic Feature Pyramid Network for UAV Object Detection}

\author{\IEEEauthorblockN{1\textsuperscript{st} Dinh Dong Tien}
\IEEEauthorblockA{\textit{Faculty of Information Technology} \\
\textit{FPT University}\\
Ho Chi Minh City, Viet Nam}
\and
\IEEEauthorblockN{2\textsuperscript{nd} Le Dang Thai An}
\IEEEauthorblockA{\textit{Faculty of Information Technology} \\
\textit{FPT University}\\
Ho Chi Minh City, Viet Nam}
\and
\IEEEauthorblockN{3\textsuperscript{rd} Phan Anh Tai}
\IEEEauthorblockA{\textit{Faculty of Information Technology} \\
\textit{FPT University}\\
Ho Chi Minh City, Viet Nam}}

\begin{document}
\maketitle

\begin{abstract}
Small-object detection in unmanned aerial vehicle (UAV) imagery is difficult
because targets occupy few pixels, appear densely, and compete with complex
backgrounds.  This paper presents Tiny-Aware Dynamic YOLO (TAD-YOLO), a
compact feature-pyramid detector built on YOLOv8n and inspired by the
shallow-oriented backbone redesign and
Tiny-focused Feature Pyramid Network (TFPN) of UAVDet.  TAD-YOLO redistributes
channel capacity toward high-resolution stages, replaces the stride-32
prediction branch with a stride-4 branch, makes the P2 output terminal, and
uses DySample in two top-down transitions.  The evaluated configuration has
1.290M parameters, requires 15.78 GFLOPs at $640\times640$, and predicts at
strides 4, 8, and 16.  A preliminary single-seed ablation on
VisDrone2019-DET shows that TAD-YOLO improves AP$_{50}$ from 31.64\% to
39.30\% and AP$_{50:95}$ from 17.42\% to 22.50\%, while reducing parameters
by 57.2\% relative to YOLOv8n.  DySample contributes +2.12 AP$_{50}$ and
+1.60 AP$_{50:95}$ over the preceding no-D5 configuration.  The results
identify P2, shallow-heavy allocation, and content-aware upsampling as the
main accuracy contributors.  Terminal P2 primarily reduces computation,
whereas D5 removal reduces both parameters and computation.
\end{abstract}

\begin{IEEEkeywords}
UAV object detection, small objects, YOLOv8, feature pyramid, P2 detection,
DySample, VisDrone.
\end{IEEEkeywords}

\section{Introduction}
Object detection from UAV platforms supports traffic monitoring, public
safety, search and rescue, and infrastructure inspection.  In contrast to
ground-level imagery, an aerial frame commonly contains many targets that are
only a few pixels wide.  Repeated spatial downsampling can erase these targets
before they reach a detection head, while background texture and severe class
imbalance increase false positives and missed detections.  The
VisDrone2019-DET benchmark makes these difficulties explicit through diverse
viewpoints, altitudes, illumination conditions, and ten traffic-related object
categories~\cite{du2019visdrone}.

YOLOv8 was selected because its implementation is strongly modular: the
backbone, C2f blocks, SPPF module, neck operations, and anchor-free decoupled
detection head are exposed as separable components~\cite{jocher2023yolov8}.
This organization permits controlled replacement of channel schedules,
upsampling operators, and prediction levels without redesigning the complete
training and inference pipeline.  The Ultralytics framework also provides a
consistent loss, augmentation, validation, checkpoint, and export interface,
which reduces implementation differences between ablations.

The nano variant, YOLOv8n, was chosen rather than a larger scale for two
reasons.  First, UAV deployment is constrained by memory, power, and latency;
a compact starting point leaves a practical budget for the expensive
high-resolution P2 branch.  Second, a lightweight baseline makes the
accuracy--complexity effect of each architectural change easier to observe
instead of hiding it inside excess model capacity.  The standard nano
configuration nevertheless allocates increasing width to deep,
low-resolution stages and predicts at strides 8, 16, and 32.  This allocation
is reasonable for general-purpose detection but may be suboptimal for the
small-object-dominated scale distribution of UAV imagery.  Retaining
high-resolution features is a direct remedy, although their cost grows
quadratically with spatial resolution~\cite{yang2022querydet}.

UAVDet addresses the same mismatch by redesigning stage depth and width and by
introducing TFPN, which replaces a large-object head with a tiny-object head
while rebalancing feature flow~\cite{yang2026uavdet}.  TAD-YOLO transfers these
principles to a different detector family: it retains the YOLOv8-style
C2f/SPPF/Detect components and does not adopt UAVDet's CSPMB or Mamba modules.
The design further introduces DySample, a point-sampling-based dynamic
upsampler~\cite{liu2023dysample}, to preserve content-dependent details during
top-down fusion.

The main contributions are as follows:
\begin{itemize}
    \item We adapt UAVDet's shallow-oriented allocation and tiny-focused
    pyramid principle to YOLOv8n through a $24/48/64/96/128$ channel schedule
    and prediction strides $\{4,8,16\}$.
    \item We formulate a terminal P2 branch: the stride-4 feature is sent
    directly to detection and is not downsampled back into the bottom-up path.
    The stride-32 D5 branch is removed.
    \item We place lightweight DySample operators in the T4-to-T3 and
    T3-to-T2 transitions while retaining nearest-neighbor interpolation for
    B5-to-T4.
    \item We conduct a component-wise preliminary ablation and explicitly
    separate measured accuracy gains from efficiency-oriented design choices.
\end{itemize}

\section{Related Work}
\subsection{Multi-scale Feature Aggregation}
Feature Pyramid Networks (FPN) introduced a top-down pathway with lateral
connections to construct semantically strong features at multiple
resolutions~\cite{lin2017fpn}.  PANet subsequently added a bottom-up path to
shorten the information route from localization-rich shallow features to deep
features~\cite{liu2018panet}.  These two ideas underlie many modern detector
necks.  TAD-YOLO retains top-down fusion but deliberately shortens bottom-up
aggregation: only the T3-to-N4 route remains.  FPN and PANet provide the
general aggregation foundation; they are not claimed as the direct source of
the terminal-P2 or no-D5 decisions.

High-resolution representations are particularly important for small
objects.  QueryDet shows that such maps improve small-object detection but
also exposes their computational cost, motivating sparse computation at
selected locations~\cite{yang2022querydet}.  TAD-YOLO takes a complementary
architectural approach.  It uses a dense P2 prediction map but compensates by
narrowing the deep stages and removing D5.

\subsection{YOLOv8 and Partial Feature Reuse}
The implementation is based on YOLOv8n~\cite{jocher2023yolov8}.  Beyond its
small parameter budget, YOLOv8n provides clean module boundaries between its
backbone, neck, and detection head.  This modularity is central to the
experimental design: the same C2f, SPPF, Detect, loss, and data pipeline are
retained while feature widths, pyramid routes, and upsampling operators are
changed.  Consequently, differences between configurations are easier to
attribute to the proposed topology rather than to a replacement detector.

The C2f blocks are conceptually related to cross-stage partial feature reuse.
CSPNet showed that separating and later merging feature paths can reduce
duplicated gradient information and computation~\cite{wang2020cspnet}.
Nevertheless, C2f is a specific Ultralytics design and should not be described
as a block proposed by CSPNet.

\subsection{UAV-Specific Small-Object Detection}
UAVDet is the central architectural inspiration for TAD-YOLO.  It combines a
redesigned CNN hierarchy, CSPMB/Mamba global--local modeling, and TFPN.  On
VisDrone, UAVDet reports improved AP and AP$_S$ with substantially fewer
parameters than its RTMDet-m baseline~\cite{yang2026uavdet}.  The inheritance
boundary is important: TAD-YOLO adopts the principles of shallow-oriented
capacity and replacing a large-object branch with a tiny-object branch, but it
does not copy the RTMDet-m backbone, CSPNeXt stages, CSPMB blocks, or Mamba
operators.  Its terminal P2 path is a project-specific realization of the
tiny-focused feature-flow principle.

SMA-YOLO provides a closely related YOLOv8n-based design point: it introduces
a $160\times160$ small-object head and removes the large-object prediction
branch~\cite{zhou2025smayolo}.  TAD-YOLO shares this small-scale emphasis but
is distinguished by shallow-heavy channel redistribution, a terminal P2
route, and DySample-based top-down refinement.  This comparison places the
proposed topology among existing tiny-object-oriented YOLO designs without
claiming that P2 or large-head removal is unique to TAD-YOLO.

\subsection{Content-Aware Upsampling}
Fixed nearest-neighbor interpolation is inexpensive but cannot adapt sampling
locations to local structure.  DySample reformulates upsampling as learned
point sampling and avoids dynamic convolution and custom CUDA
kernels~\cite{liu2023dysample}.  This combination is attractive because the
two high-resolution top-down transitions must preserve boundaries without a
large kernel-generation network.

\section{Methodology}
\subsection{Architecture Overview}
Let an RGB image be
$\mathbf{I}\in\mathbb{R}^{3\times H\times W}$, with $H$ and $W$ divisible by
32.  The backbone produces four lateral features
\begin{equation}
\mathbf{B}_i\in\mathbb{R}^{C_i\times H/2^i\times W/2^i},\quad
i\in\{2,3,4,5\},
\end{equation}
where $(C_2,C_3,C_4,C_5)=(48,64,96,128)$.  A 24-channel stride-2 stem precedes
these stages.  Their C2f repeat counts are $(1,2,2,1)$, and the B5 stage ends
with SPPF.  Compared with the conventional increasing-width pattern, this
schedule limits deep-stage width while retaining more relative capacity at
high resolution.

Figure~\ref{fig:architecture} summarizes the resulting backbone, top-down
fusion, terminal P2 output, and shortened bottom-up path.

\begin{figure*}[t]
\centering
\includesvg[width=\textwidth,inkscapelatex=false]{assets/tad_yolo_architecture}
\caption{TAD-YOLO architecture.  The shallow-heavy backbone feeds a
terminal-P2 multi-scale neck with DySample in the T4-to-T3 and T3-to-T2
transitions.  T2 is used directly by the stride-4 head and is excluded from
bottom-up aggregation; the detector predicts at strides 4, 8, and 16 without
a D5 branch.}
\label{fig:architecture}
\end{figure*}

For a $640\times640$ input, the data flow is
\begin{align}
\mathbf{T}_4 &= \mathrm{C2f}\!\left(
    U_{\mathrm{NN}}(\mathbf{B}_5)\oplus\mathbf{B}_4\right), \\
\mathbf{T}_3 &= \mathrm{C2f}\!\left(
    U_{\mathrm{Dy}}(\mathbf{T}_4)\oplus\mathbf{B}_3\right), \\
\mathbf{T}_2 &= \mathrm{C2f}\!\left(
    U_{\mathrm{Dy}}(\mathbf{T}_3)\oplus\mathbf{B}_2\right), \\
\mathbf{N}_4 &= \mathrm{C2f}\!\left(
    \mathrm{Conv}_{3\times3,s=2}(\mathbf{T}_3)\oplus\mathbf{T}_4\right).
\end{align}
Here $\oplus$ denotes channelwise concatenation, $U_{\mathrm{NN}}$ denotes
fixed $2\times$ nearest-neighbor upsampling, and $U_{\mathrm{Dy}}$ denotes
$2\times$ DySample~\cite{liu2023dysample}.  Detection is performed on
\begin{equation}
\mathcal{D}=\mathrm{Detect}\!\left(
\{\mathbf{T}_2,\mathbf{T}_3,\mathbf{N}_4\}\right),
\end{equation}
where braces denote the list of multi-scale inputs rather than a spatial or
channel concatenation.  Their strides are $(4,8,16)$ and channel counts are
$(48,64,96)$.

Table~\ref{tab:feature_contract} lists the corresponding feature dimensions
and routing roles.

\begin{table}[t]
\caption{Feature contract for a $640\times640$ input.}
\label{tab:feature_contract}
\centering
\small
\begin{tabular}{lccc}
\toprule
Feature & Channels & Resolution & Role \\
\midrule
B2 & 48  & $160\times160$ & shallow lateral \\
B3 & 64  & $80\times80$   & lateral \\
B4 & 96  & $40\times40$   & lateral \\
B5 & 128 & $20\times20$   & SPPF output \\
T2/D2 & 48 & $160\times160$ & terminal prediction \\
T3/D3 & 64 & $80\times80$ & prediction/route \\
N4/D4 & 96 & $40\times40$ & prediction \\
\bottomrule
\end{tabular}
\end{table}

\subsection{Terminal P2 and Removal of D5}
A standard bottom-up pyramid can downsample P2 into an N3 path and eventually
reach N5.  TAD-YOLO instead defines
\begin{equation}
\mathbf{T}_2\rightarrow\mathbf{D}_2,
\qquad
\mathbf{T}_2\nrightarrow\mathbf{N}_3.
\end{equation}
This terminal branch prevents the expensive $160\times160$ feature from being
reprocessed in the bottom-up path, while the stride-32 D5 head is removed.
For a $640\times640$ input, D2 contains $64\times$ as many spatial locations
as D5 ($160^2$ versus $20^2$).  The design therefore shifts prediction toward
high-resolution features.  It follows UAVDet's TFPN principle of replacing
large-object emphasis with tiny-object emphasis, whereas the terminal-P2
implementation is specific to TAD-YOLO.

\subsection{DySample Top-Down Refinement}
For each top-down input $\mathbf{X}$, DySample~\cite{liu2023dysample}
predicts grouped offsets
$\Delta\mathbf{p}$ and samples around an initial regular grid
$\mathbf{p}_0$.  In simplified form,
\begin{equation}
\mathbf{Y}(\mathbf{q})=
\mathrm{grid\_sample}\!\left(
\mathbf{X},\mathbf{p}_0(\mathbf{q})+\alpha\Delta\mathbf{p}(\mathbf{q})
\right),
\end{equation}
where the implementation uses scale 2, four channel groups, ``lp'' style,
$\alpha=0.25$, bilinear sampling, border padding, and
\texttt{align\_corners=false}.  Sampling is calculated in FP32 for numerical
stability.  DySample is used only from T4 to T3 and from T3 to T2; the
B5-to-T4 transition remains nearest-neighbor interpolation.

\subsection{Training Objective}
The detector uses the YOLOv8 detection objective~\cite{jocher2023yolov8},
including distribution focal loss (DFL) for box
regression~\cite{li2020gfl}:
\begin{equation}
\mathcal{L}=\lambda_{box}\mathcal{L}_{box}
+\lambda_{cls}\mathcal{L}_{cls}
+\lambda_{dfl}\mathcal{L}_{dfl},
\end{equation}
with $(\lambda_{box},\lambda_{cls},\lambda_{dfl})=(7.5,0.5,1.5)$.
The GFL citation identifies the origin of DFL; TAD-YOLO does not claim to
adopt the complete GFL or Quality Focal Loss objective.

\section{Experiments and Results}
\subsection{Dataset and Exploratory Analysis}
The VisDrone2019-DET train, validation, and test-dev splits comprise 8,629
images: 6,471 training, 548 validation, and 1,610 test-dev images, with ten
annotated categories~\cite{du2019visdrone}.
The converted training labels contain 344,201 instances and exhibit a
substantial long-tail distribution.  Cars
account for 42.1\% of instances and pedestrians for 23.0\%; together they form
65.1\% of all annotations.  Cars are approximately 45 times more frequent
than awning-tricycles.  Under the COCO area
thresholds, a published audit of VisDrone reports 60.49\% small objects
($a<32^2$), 33.97\% medium objects ($32^2\leq a<96^2$), and 5.53\% large
objects ($a\geq96^2$), where $a$ is the box area in square
pixels~\cite{zhou2025smayolo}.
Figure~\ref{fig:dataset_distribution} summarizes both distributions.  These
observations motivate high-resolution detection and cautious interpretation
of aggregate AP, which can hide rare-class failure modes.

\begin{figure*}[t]
\centering
\includesvg[width=\textwidth,inkscapelatex=false]{assets/visdrone_dataset_distribution}
\caption{VisDrone dataset statistics.  The bar chart shows class counts in
344,201 converted training labels; the pie chart reports the published
COCO-style object-scale proportions~\cite{zhou2025smayolo}.}
\label{fig:dataset_distribution}
\end{figure*}

\subsection{Experimental Protocol}
All recorded runs use $640\times640$ letterboxed RGB inputs and the official
training/validation split.  The common schedule is 150 epochs, batch size 16,
the hybrid Muon--SGD optimizer (MuSGD), initial learning rate 0.01, final
learning-rate factor 0.01, momentum 0.9, weight decay $5\times10^{-4}$, and
three warm-up epochs.
Augmentation includes HSV jitter, translation 0.1, scale 0.5, horizontal flip
probability 0.5, and mosaic probability 1.0; mosaic is disabled for the final
ten epochs.  The preliminary record uses seed 0 with deterministic mode enabled,
BF16 precision, PyTorch 2.10.0, and Ultralytics 8.4.154.

Per-class precision $P_c$ and recall $R_c$ are computed from true positives
(TP), false positives (FP), and false negatives (FN):
\begin{equation}
P_c=\frac{\mathrm{TP}_c}{\mathrm{TP}_c+\mathrm{FP}_c},\qquad
R_c=\frac{\mathrm{TP}_c}{\mathrm{TP}_c+\mathrm{FN}_c}.
\end{equation}
The reported $P$ and $R$ are class means.  AP$_{50}$ is mean AP across classes
at IoU 0.50, while AP$_{50:95}$ additionally averages over IoU thresholds from
0.50 to 0.95 in steps of 0.05, following the COCO-style
protocol~\cite{lin2014coco}.  AP$_{50:95}$ is the checkpoint-selection metric
because it is more sensitive to localization quality.

Parameter counts refer to trainable detector parameters.  GFLOPs are profiled
at batch size one and $640\times640$ in FP32 using $2\times$ MACs, excluding
NMS.  Functional sampling operators may be incompletely counted by THOP, so
these figures should be interpreted as architecture profiles rather than
hardware latency.

\subsection{Preliminary Ablation Results}
Table~\ref{tab:ablation} reports a sequential, single-seed ablation in which
each row explicitly names all accumulated modifications.  Adding a P2 head to
YOLOv8n produces the largest isolated gain: +4.23 AP$_{50}$ and +2.96
AP$_{50:95}$.  Adding the shallow-heavy channel schedule contributes another
+2.35 and +1.36 points, respectively, while reducing parameters from 2.928M
to 1.633M.  These changes support the hypothesis that high-resolution
capacity is valuable on VisDrone.

The results also contradict an overly simple claim that every proposed change
is independently beneficial.  Relative to YOLOv8n + P2 + shallow-heavy,
terminal P2 loses 0.58 AP$_{50}$ and 0.31 AP$_{50:95}$; removing D5 loses a
further 0.46 and 0.53 points.  Their value is primarily efficiency: the full
no-D5 predecessor retains a +5.54/+3.48 point advantage over YOLOv8n with
60.2\% fewer parameters, although GFLOPs increase from 8.20 to 13.29 because
D2 is spatially expensive.  Adding DySample then raises precision and recall
by 1.29 and 2.01 points, respectively, and improves AP$_{50}$/AP$_{50:95}$ by
2.12/1.60 points over the no-D5 predecessor.  This gain costs 0.091M
parameters and 2.49 GFLOPs.  The final TAD-YOLO configuration therefore
reaches 39.30\% AP$_{50}$ and 22.50\% AP$_{50:95}$ with 1.290M parameters.

\begin{table*}[t]
\caption{Single-seed VisDrone validation ablation.  Each row accumulates the
explicitly named modification; all accuracy values are percentages.}
\label{tab:ablation}
\centering
\small
\begin{tabular}{p{7.2cm}rrrrrr}
\toprule
Configuration & $P$ & $R$ & AP$_{50}$ & AP$_{50:95}$ & Params (M) & GFLOPs \\
\midrule
YOLOv8n & 44.16 & 32.67 & 31.64 & 17.42 & 3.013 & 8.20 \\
YOLOv8n + P2 & 47.71 & 35.93 & 35.87 & 20.38 & 2.928 & 12.37 \\
YOLOv8n + P2 + shallow-heavy & 48.20 & 38.53 & 38.22 & 21.74 & 1.633 & 14.23 \\
YOLOv8n + P2 + shallow-heavy + terminal P2 & 49.04 & 37.37 & 37.64 & 21.43 & 1.634 & 13.64 \\
YOLOv8n + P2 + shallow-heavy + terminal P2 $-$ D5 & 48.51 & 37.09 & 37.18 & 20.90 & 1.199 & 13.29 \\
\midrule
\textbf{YOLOv8n + P2 + shallow-heavy + terminal P2 $-$ D5 + DySample
(TAD-YOLO)} & \textbf{49.80} & \textbf{39.10} & \textbf{39.30} &
\textbf{22.50} & \textbf{1.290} & \textbf{15.78} \\
\bottomrule
\end{tabular}
\end{table*}

\subsection{Discussion}
P2 provides the largest initial accuracy gain, while shallow-heavy allocation
further improves AP and reduces parameters.  Terminal P2 saves 0.59 GFLOPs
with nearly unchanged parameters, and D5 removal saves a further 0.435M
parameters and 0.35 GFLOPs; both slightly reduce AP.  They therefore represent
an accuracy--efficiency trade-off rather than accuracy-enhancing modules.
DySample recovers this loss and exceeds the shallow-heavy configuration by
1.08 AP$_{50}$ and 0.76
AP$_{50:95}$, supporting content-aware high-resolution fusion.

The dense D2 head still increases computation and memory, and removing D5 may
reduce large-object recall.  Moreover, the results cover one split, resolution,
and seed; they indicate trends rather than statistical superiority.  Future
evaluation should report three-seed mean and standard deviation, per-scale and
per-class AP, peak memory, and warmed-up latency on specified hardware.

\section{Conclusion}
TAD-YOLO adapts UAVDet's shallow-oriented and tiny-focused principles to
YOLOv8n through shallow-heavy allocation, terminal P2, D5 removal, and two
DySample transitions.  It reaches 39.30\% AP$_{50}$ and 22.50\%
AP$_{50:95}$ with 1.290M parameters: gains of 7.66 and 5.08 points over
YOLOv8n with 57.2\% fewer parameters.  Multi-seed and hardware-level tests are
required to confirm robustness and deployment efficiency.

\bibliographystyle{IEEEtran}
\bibliography{references}

\end{document}
